\section{Introduction}
\label{sec:introduction}

\subsection{The Convergence Question}

Redistricting ensembles have become central to partisan-gerrymandering
litigation in state courts following \textit{Rucho v.\ Common Cause}~\citep{rucho2019},
which held that federal courts lack jurisdiction to police partisan
gerrymandering.
In the resulting state-court docket, expert witnesses routinely deploy
Markov chain Monte Carlo (MCMC) ensembles---most commonly the \recom{}
chain of \citet{deford2021}---to show that an enacted map is an outlier
relative to a distribution of algorithmically drawn plans.

The inference is straightforward: if the enacted map's Democratic seat
count falls in the 99th percentile of an ensemble of 50{,}000 plans, the
map is likely gerrymandered in favor of Republicans.
The inference is valid, however, only if the ensemble has \emph{mixed}---if
the chain has explored the relevant region of plan space thoroughly enough
that the empirical distribution of outcomes approximates the true stationary
distribution.

In practice, convergence is rarely certified.
Papers and reports describe running ``50{,}000 plans'' or ``100{,}000
steps'' without computing the diagnostics that would tell a court whether
those plans were drawn from a converged chain.
A non-converged ensemble may exhibit systematic biases---for instance,
anchoring near the starting plan---that could distort the comparison.

\subsection{Three Diagnostics}

This paper formalises three complementary convergence diagnostics for
redistricting ensembles.

\paragraph{R-hat (Potential Scale Reduction Factor).}
The Gelman-Rubin $\rhat$ statistic~\citep{gelman1992} runs multiple
independent chains from different starting plans and compares between-chain
variance to within-chain variance.
Values of $\rhat < 1.1$ indicate approximate convergence.
We extend $\rhat$ to the redistricting setting by applying it to ten
parallel \recom{} chains with randomised start plans.

\paragraph{Effective Sample Size.}
The \ess{}~\citep{geyer1992} accounts for autocorrelation in the chain.
An MCMC chain with high autocorrelation produces a smaller effective
sample than its nominal step count suggests.
For redistricting chains with typical lag-1 Hamming autocorrelation
$\hat{\rho}_1 \approx 0.85$--$0.95$, a 10{,}000-step chain yields
$\ess \approx 500$--$2{,}000$ for Polsby-Popper compactness.

\paragraph{Hamming Autocorrelation.}
The Hamming distance between two redistricting plans is the fraction of
census tracts assigned to different districts.
The lag-$k$ Hamming autocorrelation $\ham(k)$ measures how rapidly
consecutive plans decorrelate.
For \recom{}, typical values are $\ham(1) \approx 0.85$--$0.95$ and
$\ham(k) \approx \hat{\rho}_1^k$ (geometric decay).

\subsection{Main Results}

We apply all three diagnostics to ten parallel \recom{} chains across six
states: North Carolina (14 districts), Wisconsin (8), Georgia (14),
Pennsylvania (17), Texas (38), and California (52).
Section~\ref{sec:application} presents a table of $\rhat$, $\ess$, and
$\ham(1)$ for each state.
Key findings:

\begin{enumerate}
\item $\rhat < 1.1$ is achieved after 2{,}000--20{,}000 steps depending
      on the state; California (large, compact) converges most slowly.
\item Typical $\ess$ at 10{,}000 steps is 500--2{,}000 for \PP{} and
      300--1{,}000 for Democratic seat counts.
\item $\ham(1) \approx 0.87$--$0.93$ is consistent across states and
      largely independent of $k$.
\end{enumerate}

\subsection{The ConvergenceSweep Alternative}

The \cs{} algorithm~\citep{b16paper} provides a deterministic, $O(T \cdot n)$
stopping criterion: sweep forward from the content-derived seed until
$T = 600$ consecutive seeds produce no improvement.
This is not an MCMC method; it does not sample from a distribution.
For the \dia{}'s plan-\emph{generation} task, \cs{} is preferable because
it produces a certifiably optimal plan without relying on asymptotic
convergence theory.
For the distinct task of characterising the feasible plan space and
auditing enacted maps, ensemble methods remain irreplaceable.
Section~\ref{sec:statutory} formalises this division of labor.

\subsection{Implementation}

All three diagnostics are implemented in the \texttt{redist} Rust binary:
\begin{verbatim}
  redist analyze --types ensemble-diagnostics --state NC --year 2020
\end{verbatim}
The command runs ten chains, computes $\rhat$, $\ess$, and $\ham(1)$,
and reports a pass/fail verdict against the thresholds defined in
Section~\ref{sec:statutory}.
